\section{B. Adelic approximations, completions, and exactness}
\subsection{The definition actually given in \S52}

\begin{statement}{B.1}{Proved}{The axioms of a Galois--Teichm\"uller structure}
The final provisional definition in \MS{52}{750--752}{205--209}
consists of inclusions
\[
 \Pi_0\subset\mathcal G_0,\quad \Pi\subset\mathcal G\subset\mathcal M,
 \qquad \Pi=\widehat{\Pi_0},\quad \mathcal G=\overline{\mathcal G_0},
\]
with the second row profinite, the first discrete, $\Pi,\mathcal G$ normal
in $\mathcal M$, the specified peripheral $L$-structure invariant under
$\mathcal M$, and $\mathcal G_0$ normalizing $\Pi_0$. Its cyclotomic
character factors through $\Gamma=\mathcal M/\mathcal G$ and has open
image in $\widehat\Z^\times$. We do not require
$\widehat{\mathcal G_0}\cong\mathcal G$, which the manuscript expressly
does not require at scan206.

A Lie approximation is an affine finite-type $\Q$-group $G$, a character
$\chi_G:G\to\Gm$, and a continuous homomorphism
$\rho:\mathcal M\to G(\mathbb A_f)$ such that
\[
 \chi_G\rho=\chi,\qquad \rho(\mathcal G_0)\subset G(\Q),
\]
where $\Q$ is embedded diagonally, and such that no proper closed
$\Q$-subgroup contains $\rho(\mathcal M)$ adelically.
These are conditions $1^\circ,2^\circ,3^\circ$ forte, and $4^\circ$
of scans202--208.
\end{statement}
\begin{proof}
Diagram (35), visually checked at scan208, factors through $G(\Q)$;
the ambient representation has target $G(\mathcal A_\Q)$, not
$G(\Q_\ell)$ for a single prime. The strengthening from $\Pi_0$
to $\mathcal G_0$ is explicit in scans204--208.
The arithmetic fundamental-group sequence of a hyperbolic curve over a
finitely generated field of characteristic zero supplies the case
$\mathcal G=\Pi$: choose a complex comparison and its discrete
peripheral group. For the universal smooth curve, comparison with
the complex topological fibre and moduli stack supplies the
geometric groups and their closures. The étale homotopy sequence
gives normality; conjugation on a puncture's inertia is cyclotomic,
because an $n$th-root local parameter is multiplied by an $n$th
root of unity. The character is trivial geometrically. The algebraic
constants in the base field have finite degree over $\Q$, hence
their intersection with $\Q(\mu_\infty)$ is finite; the cyclotomic
image is open. These checks do not assert the congruence-subgroup
property for mapping class groups. With no cusps a cyclotomic
character is extra data, not recoverable from an empty list.
\end{proof}

\begin{statement}{B.2}{Proved}{A universal pro-object of admissible targets}
For fixed B.1 data, let morphisms of admissible targets be
$\Q$-group morphisms compatible with $\rho$ and $\chi_G$.
This category is cofiltered and has at most one morphism between
two objects. Its inverse limit $\mathscr E$ is a pro-affine
$\Q$-group universal for the specified admissible finite-type
targets. This is not an identification with the absolute
proalgebraic envelope of $\Pi_0$ or its Malcev completion.
\end{statement}
\begin{proof}
The object $\Gm$ with $\rho=\chi$ exists. Given two targets take
the scheme-theoretic $\Q$-Zariski closure of the image in their
product over $\Gm$. Intersecting the closed subgroups containing
the image gives a closed subgroup; finite type follows from
noetherianity of the ambient coordinate ring. This is a common
refinement. Continuity, rationality on $\mathcal G_0$, and the
character condition are preserved, and minimality is built into
the closure. Two morphisms agreeing on the adelic image agree on
its Zariski closure and hence everywhere. Choose a set of
representatives of the finite-type coordinate algebras and
representations. The direct limit of their coordinate Hopf
algebras represents the claimed inverse limit. This proves its
universal property without imposing unipotence or weights.
\end{proof}

\begin{statement}{B.3}{Refuted}{Lead 1 as an identification with a nilpotent tower}
B.1 is not the replacement of the geometric groups by nilpotent
or solvable completions. It specifies neither a first nor a second
stage. A pair of first two approximations therefore requires a
chosen filtration or subcategory of targets.
\end{statement}
\begin{proof}
For the elliptic modular arithmetic fundamental group take the
adelic Tate-module representation in $\GL_{2,\Q}$, with determinant
as character. Topological comparison sends the geometric discrete
group $\SL_2(\Z)$ to its usual rational matrices. Its Zariski
closure is $\SL_{2,\Q}$: it contains both upper and lower
one-parameter unipotent subgroups, whose brackets generate
$\mathfrak{sl}_2$. The determinant has open cyclotomic image,
so the full closure is $\GL_{2,\Q}$. This is an admissible target
with non-solvable connected geometric group. A construction
making the whole geometric group prounipotent cannot have this
image. Weighted completion, relative completion with a fixed
reductive representation, and genus-zero or elliptic $\GT$
specify different universal problems; the label Lie approximation
alone supplies no comparison map. The strongest corrected
interpretation is the category of B.1 and its pro-object B.2.
\end{proof}

\begin{statement}{B.4}{Proved}{Two explicitly chosen Malcev stages}
For types $(0,3),(0,4),(1,1)$ the geometric discrete group is
free of rank $r=2,3,2$. Choose rational Malcev nilpotence-class
truncation. Its first stages are
\[
 U_1=\Ga^r,\qquad
 U_2=\bigl(V\oplus\Lambda^2V,*\bigr),\qquad
 (v,a)*(w,b)=(v+w,a+b+\tfrac12v\wedge w).
\]
The Lie bracket is $[v,w]=v\wedge w$ and $\Lambda^2V$ is central;
lower-central degrees are $1,2$. For punctured spheres the graded
Galois representations are $\Q_\ell(1)^r$ and
$\Q_\ell(2)^{\binom r2}$; for a punctured elliptic curve they are
$T_\ell E\otimes\Q_\ell$ and $\Q_\ell(1)$. Permutations of punctures
act on $V$ by permuting peripheral classes subject to their sum
being zero, and act by the exterior square in degree two.
\end{statement}
\begin{proof}
The displayed Lie algebra has the universal property of the
free class-two Lie algebra: arbitrary images of a basis of $V$
extend uniquely, and all brackets of length three vanish.
Truncated Baker--Campbell--Hausdorff gives the group law;
bilinearity verifies associativity and inverses are $(-v,-a)$.
Cutting the three surfaces along standard arcs gives their free
fundamental groups. Functoriality determines the graded action
from abelianization. Peripheral inertia is $\Z_\ell(1)$; for the
torus the Weil pairing identifies $\Lambda^2T_\ell E$ with
$\Z_\ell(1)$. This determines the graded action, not a canonical
splitting of a filtered genuine action. Base-point cocycles and
integral models are additional choices (Parts D, E, G); in
particular the factor $1/2$ requires attention at $2$.
\end{proof}

\subsection{A covering-exact family within the scope of \S50}

\begin{statement}{B.5}{Proved}{Virtually abelian quotients and finite covers}
For a finitely generated discrete group $F$ put
\[
 V(F)=\varprojlim_{K\triangleleft F,\ [F:K]<\infty}F/[K,K].
\]
If $F\twoheadrightarrow H$ is a finite quotient with kernel $N$,
then $1\to V(N)\to V(F)\to H\to1$ is exact as an inverse limit
of discrete groups. For free $F$ the \S50 envelopes apply to
every stage, and their inverse system is also exact for this
finite covering.
\end{statement}
\begin{proof}
Restrict to $K\subset N$, a cofinal subsystem. It is also cofinal
among the finite-index normal subgroups of $N$: intersect finitely
many $F$-conjugates of such a subgroup. At each stage,
$1\to N/[K,K]\to F/[K,K]\to H\to1$ is exact. Kernels commute
with inverse limits, and each element of $H$ has compatible lifts
given by a fixed representative in $F$.

If $F$ is free, $K^{\ab}$ is a finite-rank lattice. The two groups
$N/[K,K]\subset F/[K,K]$ have this same normal lattice.
In the explicit \S50 construction, each finite scheme stage is
a disjoint union of affine spaces with multiplication determined
by an integral extension cocycle. The subgroup over $N/K$ is
the union of the corresponding components, hence its envelope
is exactly the fibre over the identity of the constant group
$H$. Transition maps preserve these component decompositions;
inverse limits prove the assertion for the envelopes. Without
the lattice hypothesis only the group inverse-limit statement
has been asserted here.
\end{proof}

\begin{statement}{B.6}{Proved}{Fundamental localizers of the specified completions}
For abelianization, the metabelian quotient, a fixed positive
nilpotence-class quotient, pro-$\ell$ completion, and rational
prounipotent completion, the class $W_{\mathfrak a}$ in the
prompt satisfies two-out-of-three, retract closure, the
terminal-object axiom, and relative Theorem A. The literal
\S52 construction requires richer arithmetic data and does not
define a functor of an arbitrary bare fundamental group.
\end{statement}
\begin{proof}
Use all coefficient groups $H$ which are, respectively, abelian,
metabelian, nilpotent of the fixed class, finite $\ell$-groups,
or groups $U(\Q)$ for finite-dimensional unipotent $\Q$-groups.
For a connected space $X$, $\operatorname{Map}(X,BH)$ is the
nerve of the groupoid of homomorphisms $\pi_1X\to H$ and their
conjugations. The completion's universal property identifies
these representations, and also their centralizers: finite
images are unchanged; in the unipotent case the Zariski closures
of the images are unchanged. Conversely equivalence on all
these mapping spaces determines the quotient or pro-object by
its universal property. Including discrete sets as targets
detects $\pi_0$. Thus $W_{\mathfrak a}$ is exactly the class
inverted by all these mapping-space functors.

Two-out-of-three and retract closure follow. Mapping out of a
homotopy colimit gives a homotopy limit. For a triangle
$A\to B\to S$ of small categories, the Grothendieck construction
gives
\[
 BA\simeq\operatorname*{hocolim}_{s\in S}B(A/s),\qquad
 BB\simeq\operatorname*{hocolim}_{s\in S}B(B/s).
\]
The first comparison follows by projecting the category of
pairs $(a,f(a)\to s)$ to $A$: the arrows out of $f(a)$ form
a category with initial object its identity. The second is
identical. An objectwise equivalence after mapping to every
test target therefore gives the required global equivalence.
This proves relative Theorem A. A terminal object contracts the
nerve, proving its axiom. This argument does not imply that all
coverings are intrinsic fibrations for these localizers.
\end{proof}

\begin{statement}{B.7}{Proved}{Typed comparisons for finite quotients}
Let $F=F_2$, $F\twoheadrightarrow H$, $d=|H|$, and $N$ its kernel.
The prompt's comparison map exists only when this finite quotient
factors through the chosen approximation. In that case it is
surjective, with the following images and kernels:
\[
\begin{array}{c|c|c}
\text{approximation}&\text{image}&\text{kernel}\\ \hline
(-)^{\ab}&N/F'&(N\cap F')/N'\\
(-)/(-)''&N/F''&F''/N''\\
(-)/\gamma_{c+1}&N/\gamma_{c+1}F&
                    \gamma_{c+1}F/\gamma_{c+1}N .
\end{array}
\]
The required conditions are respectively $H$ abelian, $H''=1$,
and $\gamma_{c+1}H=1$. The first kernel is free abelian of rank
$d-1$. None of these three maps is an isomorphism for $d>1$.
Pro-$\ell$ completion is exact for all finite $\ell$-group
quotients; other finite groups cannot be its continuous quotients.
A connected prounipotent $\Q$-group has no nontrivial finite
étale quotient.
\end{statement}
\begin{proof}
The quotient formulas are the third isomorphism theorem.
Schreier gives $\rank N=d+1$. In the abelian row its image is
a rank-two lattice. In the nilpotent row the target has
rational abelianization of rank two, since its Malcev Lie algebra
is the free class-$c$ Lie algebra on two generators, whereas the
source has abelianization of rank $d+1$.

For the metabelian row choose $v\in N\cap F'$ nonzero in
$N^{\ab}$, using the same rank inequality. The Magnus matrix
embedding identifies the derived subgroup of a free metabelian
group with a submodule of a free module over the integral domain
$\Z[N^{\ab}]$. Conjugation by $v$ is multiplication by a
nonidentity monomial, so $v$ fails to centralize some
$w\in N'/N''$. Thus $[v,w]\ne1$ modulo $N''$, but
$v,w\in F'$ imply $[v,w]\in F''$. This proves a nonzero kernel.
Explicitly the Magnus embedding sends a free basis $t_i$ to
$\left(\begin{smallmatrix}T_i&e_i\\0&1\end{smallmatrix}\right)$
over the Laurent-polynomial module and has kernel $N''$.

For an $\ell$-group $H$ the induced pro-$\ell$ topology on $N$
is its full pro-$\ell$ topology: intersect the finitely many
$F$-conjugates of a normal subgroup of $\ell$-power index in
$N$. Its quotient and the quotient $F/N$ have $\ell$-power
orders, so the intersection is open in the pro-$\ell$ topology
of $F$. Completion is therefore exact. The converses follow
from the elementary finite-quotient properties stated. For
unipotent groups, connectedness forces every map to a finite
constant group to be trivial at each finite-dimensional stage.
The missing absolute-unipotent map is not an inexact map.
\end{proof}

\begin{statement}{B.8}{Proved}{The Fermat cover and relative completion}
For $N=\ker(F_2\to(\Z/n)^2)$, $n>1$, the absolute map
$N^{\un}_\Q\to F_2^{\un}_\Q$ is surjective. On abelianizations
its source, target, and kernel have dimensions $n^2+1$, $2$,
and $n^2-1$. Relative completion with respect to
$H=(\Z/n)^2$ instead gives
\[
 1\to N^{\un}_\Q\to\mathscr R_H(F_2)\to H_\Q\to1.
\]
Its abelian unipotent kernel has dimension $n^2+1$, and
$H_1(N,\Q)\cong\Q\oplus\Q[H]$ as an $H$-module.
\end{statement}
\begin{proof}
Schreier gives $\rank N=1+n^2$. A finite-index subgroup has
Zariski-dense image in every connected unipotent quotient of
$F_2$: finitely many cosets of its image closure cover the group,
so the closure has finite index and equals the connected group.
Equivalently the map on rational $H_1$ is onto, and the free
pronilpotent Lie algebra universal property gives surjectivity.

Construct the relative group as the inverse limit of unipotent
extensions of $H_\Q$ with a lift of $F_2\to H$. Its kernel
has exactly the universal property of $N^{\un}_\Q$: induction
over the finitely many cosets embeds every unipotent
$N$-representation in a finite-component representation of
$F_2$, and restriction gives the inverse comparison of
universal properties. The cellular chains of the regular
covering of a two-petalled rose give
$0\to H_1\to\Q[H]^2\to\Q[H]\to\Q\to0$.
Semisimplicity of $\Q[H]$ proves the module formula.
Its nontrivial finite-order eigenspaces cannot occur in a
unipotent action, because a finite-order unipotent matrix in
characteristic zero is the identity. This is the precise
information forgotten by absolute completion; a group of
conjugator parameters of dimension six is a different object.
\end{proof}

\begin{statement}{B.9}{Proved}{Intrinsic finite coverings}
Finite coverings intrinsic for $W_{\mathrm{pro}\text{-}\ell}$
are precisely those with $\ell$-group monodromy. Those intrinsic
for rational unipotent completion are precisely those with
trivial monodromy. A finite covering of $BF_2$ with nontrivial
monodromy is intrinsic for neither abelianization nor any fixed
positive nilpotence-class approximation.
\end{statement}
\begin{proof}
Pro-$\ell$ equivalences preserve finite $\ell$-sets and the
completions of their stabilizers by the induced-topology
argument of B.7. This proves preservation after every base
change. If finite monodromy is not an $\ell$-group, an element
of prime order $q\ne\ell$ acts nontrivially on the fibre.
Pull back along a loop representing it. Multiplication by $q$
on $B\Z$ is a pro-$\ell$ equivalence but splits a nontrivial
orbit into singletons, failing the component test. For rational
unipotent completion the same proof uses the order of any
nontrivial monodromy element.

For the last assertion pull back to a two-petalled rose by
sending $x$ to a nontrivial monodromy loop and $y$ to the
identity. One orbit gives a cyclic $m$-sheeted cover, $m>1$.
The endomorphism $\phi(x)=x$, $\phi(y)=yx^{-1}y^{-1}xy$ is the
identity on abelianization and an automorphism on every
nilpotent quotient: it is the identity on its associated
graded, so induction along the central extensions applies.
It preserves cyclic monodromy. A Schreier basis is
$a=x^m$, $b_i=x^iyx^{-i}$, $0\le i<m$. On its abelianization,
$a$ is fixed and $b_i\mapsto2b_i-b_{i-1}$, with cyclic
subscripts. Here the group commutator used in the prompt is $[x,y]=x^{-1}y^{-1}xy$.
The determinant is $2^m-1\ne\pm1$. Thus the pullback is
not an abelian equivalence and not a fixed-class nilpotent
equivalence. This also separates exactness of a comparison on
one cover from universal preservation of equivalences.
\end{proof}
